\documentclass[11pt]{article}

\usepackage[margin=1in]{geometry}
\usepackage{amsmath,amssymb}
\usepackage{graphicx}
\usepackage{natbib}

\title{Predictive e-diagnostics for multi-state movement models}
\author{Author information to be completed}
\date{Draft version: 2026-05-07}

\begin{document}
\maketitle

\begin{abstract}
Hidden Markov models are widely used to represent behavioural state switching in animal movement data. This manuscript develops a predictive diagnostic framework in which a fitted movement HMM is evaluated through its sequential observable predictive distribution. Diagnostic alternatives are encoded as predictive density ratios, producing e-values and e-processes that can be monitored over time. The intended contribution is a statistically controlled and ecologically interpretable workflow for detecting and localizing generative model failures.
\end{abstract}

\section{Introduction}

Hidden Markov models provide a flexible framework for modelling animal movement as a state-switching process. Existing R tools such as \citet{michelot_2016_movehmm} and \citet{mcclintock_2018_momentuhmm} support practical analyses of step lengths, turning angles and covariate-dependent movement processes. This paper studies how such models can be diagnosed as sequential generators using e-values and e-processes \citep{vovk_wang_2021_evalues,howard_2020_time_uniform}.

\section{Background}

\subsection{Multi-state movement models}

To be completed.

\subsection{Predictive distributions in HMMs}

To be completed.

\subsection{E-values and e-processes}

To be completed.

\section{Predictive e-diagnostics}

\subsection{Sequential prediction under the null HMM}

To be completed.

\subsection{Diagnostic alternatives}

To be completed.

\subsection{Anytime-valid monitoring}

To be completed.

\section{Movement-specific diagnostics}

To be completed.

\section{Simulation study}

To be completed.

\section{Application}

To be completed.

\section{Discussion}

To be completed.

\bibliographystyle{plainnat}
\bibliography{references}

\end{document}
